%% =====================================================================
%% PWM-PET-IQ 1.0 Dataset Paper --- Nature Scientific Data submission
%%
%% Target venue: Nature Scientific Data (Data Descriptor format)
%%
%% Strategy: This descriptor is the DATA-DESCRIPTOR sibling of the WS-2
%% signal-equivalence framework paper. Where WS-2 (Nature Methods target)
%% contributes a *method*, this paper contributes a *dataset*: paired
%% full- and reduced-count NEMA NU-2 IQ phantom PET acquisitions for
%% reduced-activity reconstruction benchmarking. The pwm_dose_equivalence
%% framework appears here only as a Usage-Notes consumer of the dataset,
%% not as the headline object (Scientific Data does not publish methods).
%%
%% Status: 1.0 SKELETON. Every \todo{} is an acquisition- or deposit-gated
%%         blank; the dataset itself must be acquired and deposited before
%%         submission. NO measurement in this draft is asserted from a
%%         committed artifact yet --- see SUBMISSION_CHECKLIST.md, which
%%         maps each \todo{} to the acquisition/analysis/deposit step that
%%         fills it.
%%
%% Build:  pdflatex manuscript && bibtex manuscript && pdflatex manuscript
%% =====================================================================

\documentclass[twoside,11pt]{article}

%% --- Packages ---------------------------------------------------------
\usepackage[utf8]{inputenc}
\usepackage[T1]{fontenc}
\usepackage{lmodern}
\usepackage{microtype}
\usepackage[margin=1in]{geometry}
\usepackage{graphicx}
\usepackage{booktabs}
\usepackage{tabularx}
\usepackage{multirow}
\usepackage{amsmath,amssymb}
\usepackage{siunitx}
\usepackage{xcolor}
\usepackage[colorlinks=true,linkcolor=blue,citecolor=blue,urlcolor=blue]{hyperref}
\usepackage{xurl}
\usepackage[round]{natbib}
\usepackage{authblk}
\usepackage{lineno}
\linenumbers

%% --- TODO macros ------------------------------------------------------
\newcommand{\todo}[1]{\textcolor{red}{\textbf{[TODO: #1]}}}
\newcommand{\todoTable}[2]{%
  \begin{table}[ht]\centering%
    \fbox{\parbox{0.85\linewidth}{\centering\todo{TABLE: #2}}}%
    \caption{#2 (placeholder; populate when data lands).}\label{#1}%
  \end{table}}
\newcommand{\todoFigure}[2]{%
  \begin{figure}[ht]\centering%
    \fbox{\parbox{0.85\linewidth}{\centering\todo{FIGURE: #2}}}%
    \caption{#2 (placeholder; populate when data lands).}\label{#1}%
  \end{figure}}

%% --- Title block ------------------------------------------------------
\title{\textbf{PWM-PET-IQ 1.0: Paired Full- and Reduced-Count NEMA NU-2 Image-Quality Phantom PET Acquisitions for Reduced-Activity Reconstruction Benchmarking}}

\author[1,2]{\textit{Authors to be confirmed at submission}}
\affil[1]{University of Texas Southwestern Medical Center, Dallas, TX, USA}
\affil[2]{PWM Protocol Foundation}

\date{Release 1.0 --- \today}

\begin{document}
\maketitle

%% =====================================================================
\begin{abstract}
\noindent Reduced-activity (``low-count'') PET reconstruction methods are
proliferating, but they are benchmarked on heterogeneous, often
patient-specific data with no physically-known ground truth and no shared
reduced-count reference, so their reported gains are not comparable across
studies. We present \textbf{PWM-PET-IQ 1.0}, a phantom PET dataset designed
specifically as a shared, ground-truth-anchored substrate for reduced-activity
reconstruction benchmarking. The dataset comprises list-mode PET acquisitions
of the NEMA NU-2 image-quality phantom (six fillable spheres of known diameter
plus a low-density lung insert) at a sphere-to-background ratio of
\todo{SBR, e.g. 4:1}, together with (i) genuinely reduced-activity acquisitions
of the identical phantom fill obtained across radiotracer decay at
\todo{K} nominal count levels ($r \in \{1.0, 0.50, 0.25, 0.10\}$ of the
full-count reference) and (ii) \todo{N} controlled reduced-count realizations
per level produced by seeded list-mode Poisson thinning of the full-count
acquisition, so that both physically-decayed and statistically-thinned
reduced-count data are available for the same phantom geometry. All series are
reconstructed under a single fixed protocol and released as list-mode (where
license permits), sinograms, and reconstructed volumes with harmonized
per-series metadata. Because the phantom's true activity concentrations and
sphere sizes are known by construction, every image-quality endpoint
(contrast-recovery coefficient per sphere, background variability, contrast-to-noise,
and the discrepancy between decayed and thinned reduced-count data) is validated
against physical ground truth rather than against a reader panel. The release is
identified by a SHA-256 content manifest and deposited under a PhysioNet DOI with
a mirrored Zenodo archive. As a phantom acquisition it contains no protected
health information. A pip-installable loader and the companion
\texttt{pwm\_dose\_equivalence} tool consume the dataset directly (Usage Notes).
\end{abstract}

%% =====================================================================
\section*{Background \& Summary}

Reducing the injected radiotracer activity in positron-emission tomography lowers
patient radiation exposure and enables faster or cheaper imaging, at the cost of
higher image noise and reduced lesion detectability. A large and growing body of
reconstruction and denoising methods --- iterative, model-based, and
deep-learning --- targets this reduced-activity regime~\citep{lowcountpet,reader2021artificial}.
Yet these methods are difficult to compare across papers for three reasons that a
phantom dataset is uniquely positioned to remove.

\begin{enumerate}
    \item \textbf{No physical ground truth.} On patient data, the ``true'' lesion
    contrast is unknown, so reduced-activity image quality can only be scored
    against a full-dose \emph{reconstruction}, not against truth. A method that
    reproduces the full-dose reconstruction's own errors scores well. The NEMA
    NU-2 image-quality phantom has spheres of \emph{known} diameter and a fill of
    \emph{known} activity concentration, so contrast recovery is measured against
    physical truth.
    \item \textbf{No shared reduced-count reference.} Each group generates its own
    low-count data --- by injecting less activity, by shortening the frame, or by
    thinning list-mode counts --- and these are not interchangeable. There is no
    public released dataset that provides \emph{both} a genuinely reduced-activity
    acquisition and a controlled statistically-thinned realization of the same
    phantom, so the community cannot even agree on what ``25\% activity'' data
    should look like.
    \item \textbf{No count-matched noise realizations.} Robust evaluation of a
    denoiser requires multiple independent noise realizations at each count level;
    single-acquisition phantom releases provide one realization per level and
    cannot support variance estimation.
\end{enumerate}

\textbf{PWM-PET-IQ 1.0} is built to close these three gaps. Its
\emph{new, openly-licensed data records} are (a) a high-count reference list-mode
acquisition of the NEMA NU-2 IQ phantom; (b) genuinely reduced-activity
acquisitions of the identical phantom fill acquired across radiotracer decay at
\todo{K} nominal levels; (c) \todo{N} seeded list-mode-Poisson-thinned reduced-count
realizations per level, giving controlled, reproducible noise replicates; (d)
single-protocol reconstructions of every series; and (e) harmonized per-series
metadata, a physically-known ground-truth table, deterministic content addressing,
and a loader. Because both physically-decayed and statistically-thinned
reduced-count data are present for the same geometry, the dataset also lets users
\emph{measure} the discrepancy between the list-mode-thinning model (the
count-reduction model assumed by most low-count PET simulators and by the
companion signal-equivalence framework~\citep{ws2framework}) and true
reduced-activity acquisition --- a quantity that is asserted rather than measured
in most of the low-count PET literature. Table~\ref{tab:series_overview} summarizes
the release.

\begin{table}[ht]
    \centering
    \small
    \setlength{\tabcolsep}{4pt}
    \begin{tabularx}{\textwidth}{@{}>{\raggedright\arraybackslash}X c c c >{\raggedright\arraybackslash}X@{}}
        \toprule
        \textbf{Series class} & \textbf{Count level(s) $r$} & \textbf{Realizations} & \textbf{Origin} & \textbf{Purpose} \\
        \midrule
        Full-count reference       & 1.0                     & 1              & Physical acquisition & Reference / ground-truth-anchored recon \\
        Decayed reduced-activity   & \todo{levels}           & 1 per level    & Physical acquisition (across decay) & True reduced-activity reference \\
        Thinned reduced-count      & \todo{levels}           & \todo{N} per level & List-mode Poisson thinning (seeded) & Controlled noise replicates \\
        \midrule
        \textbf{1.0 total series}  & ---                     & \textbf{\todo{total}} & --- & --- \\
        \bottomrule
    \end{tabularx}
    \caption{Composition of PWM-PET-IQ 1.0. The full-count reference and the
    decayed reduced-activity series are physical acquisitions of the identical
    phantom fill; the thinned reduced-count series are seeded list-mode
    Poisson-thinned realizations of the full-count reference at matched nominal
    count levels. Exact count levels, realization counts, and totals are filled
    from the deposited acquisition log at submission (SUBMISSION\_CHECKLIST \S1--\S2).}
    \label{tab:series_overview}
\end{table}

The dataset is a phantom acquisition, so it carries no protected health
information and requires no patient consent --- a clean fit for open release. The
companion signal-equivalence framework~\citep{ws2framework} and the PWM-LDCT CT
dataset~\citep{ws1dataset} are separate releases; this descriptor is
self-contained and can be obtained, used, and validated without either.

%% =====================================================================
\section*{Methods}

\subsection*{Ethics and data sharing}

PWM-PET-IQ 1.0 is a physical phantom acquisition. No human or animal subjects are
involved, no protected health information is generated, and no institutional
review-board approval or informed consent is required; the acquisition site's
radiation-safety authorization for phantom filling is on file
(\todo{radiation-safety authorization no.}). The dataset is released under
\todo{license, e.g. CC BY 4.0} with no redistribution restriction.

\subsection*{Phantom}

We used the NEMA NU-2 image-quality (IQ) body phantom~\citep{nema_nu2}, comprising
a body-cavity background compartment, a centrally-mounted low-density cylindrical
lung insert, and six fillable spheres of internal diameter 10, 13, 17, 22, 28, and
37\,mm. The four smallest spheres (10--22\,mm) are conventionally treated as
``hot'' (filled above background) and the two largest (28, 37\,mm) as either hot or
cold depending on protocol; the hot/cold assignment used here is
\todo{hot/cold sphere assignment}. Sphere and background fill activity
concentrations were prepared to a sphere-to-background ratio of \todo{SBR, e.g. 4:1},
with background activity concentration \todo{background conc. (kBq/mL)} at reference
start time. Fill activities were assayed in a dose calibrator
(\todo{dose-calibrator model}) cross-calibrated to the PET scanner within
\todo{cross-calibration window}.

\subsection*{Radiotracer and activity levels}

The phantom was filled with \todo{radiotracer, e.g. $^{18}$F-FDG}
(half-life \todo{$T_{1/2}$}). The full-count reference acquisition began at a total
phantom activity of \todo{reference activity (MBq)}. Genuinely reduced-activity
series were obtained by re-acquiring the \emph{identical} phantom fill at later
timepoints as the tracer decayed, targeting nominal count fractions
$r \in \{1.0, 0.50, 0.25, 0.10\}$ relative to the reference; the decay timepoints,
elapsed times, and decay-corrected activities are recorded per series in
\texttt{metadata.json} and summarized in Table~\ref{tab:acq_params}. Acquiring
across decay holds phantom geometry, positioning, and fill exactly constant while
varying only the count level, which is the property that makes the decayed series a
valid physical reference for the thinned series.

\subsection*{Acquisition}

All series were acquired on \todo{scanner model} in list-mode with
\todo{TOF/PSF capabilities} enabled. The full-count reference was acquired for
\todo{reference frame duration} at a single bed position centred on the spheres;
each decayed reduced-activity series was acquired for \todo{decayed frame duration}
at the identical bed position. Acquisition geometry (bed position, phantom
orientation, axial FOV) was held fixed across all physical series and is recorded
in \texttt{metadata.json}.

\subsection*{Reduced-count realization by list-mode thinning}

For each target count fraction $r$, we generated \todo{N} independent reduced-count
realizations from the full-count reference list-mode stream by binomial (Poisson)
thinning: each recorded coincidence event is retained independently with
probability $r$, using a counter-based deterministic RNG seeded by
$(\texttt{series\_id}, r, \texttt{realization\_index})$ so that every realization is
bit-reproducible from the released full-count list-mode and the published seed
recipe. Thinning at rate $r$ reproduces the count statistics of an acquisition at
$r\times$ the injected activity (same scan duration) at the level of the marginal
coincidence-count distribution; it does \emph{not} reproduce activity-dependent
detector effects (dead-time, randoms fraction, pile-up), which differ between a
thinned stream and a truly lower-activity acquisition. The magnitude of that
discrepancy is exactly what the paired decayed-vs-thinned comparison in Technical
Validation quantifies, rather than assuming it away.

\subsection*{Reconstruction}

Every series --- reference, decayed, and thinned --- was reconstructed under a
single fixed protocol to eliminate reconstruction as a confound:
\todo{algorithm, e.g. OSEM}~\citep{hudson1994osem} with \todo{iterations}
iterations $\times$ \todo{subsets} subsets, \todo{matrix size} matrix,
\todo{voxel size (mm)} voxels, \todo{PSF on/off} resolution modelling,
\todo{TOF on/off} time-of-flight, and a \todo{post-filter (mm FWHM)} Gaussian
post-filter, with all standard corrections (attenuation from the
\todo{attenuation source, e.g. CT}, scatter, randoms, normalization, decay, and
dead-time) applied. Reconstruction settings are recorded verbatim in
\texttt{metadata.json} and pinned in the released reconstruction container
(Data Records). An additional filtered-back-projection reconstruction is provided
per series as an analytic reference where the scanner toolchain permits
(\todo{FBP availability}).

\subsection*{Ground-truth table}

Because the phantom is physical, ground truth is known by construction rather than
adjudicated. The released \texttt{ground\_truth.json} records, per sphere, the
nominal internal diameter, the measured fill activity concentration (decay-corrected
to each series' start time), and the resulting true sphere-to-background contrast;
for the background and lung-insert compartments it records the true concentration
and the geometry of the standard 12-region background ROI set specified by NEMA
NU-2~\citep{nema_nu2}. All Technical-Validation endpoints are computed against this
table.

\subsection*{De-identification and provenance}

Phantom acquisitions contain no PHI; nonetheless every released DICOM object is
passed through the same tag-whitelisting pipeline used for the CT
release~\citep{ws1dataset} to strip operator names, institution strings, and free-text
fields, and the strip log is deposited as Supplementary Material. Each released
artifact is recorded in a SHA-256 content manifest (Data Records).

%% =====================================================================
\section*{Data Records}

PWM-PET-IQ 1.0 is deposited under a PhysioNet DOI (\todo{PhysioNet URL}) and
mirrored to Zenodo (\todo{Zenodo DOI}); the identical content is addressed by the
SHA-256 manifest \texttt{manifest.sha256}. Unlike the CT release, the phantom scans
carry no DUA restriction, so \emph{all} pixel and list-mode records are deposited in
full (subject only to the list-mode-format license note below). The release layout is
given in Table~\ref{tab:data_records}.

\begin{table}[ht]
    \centering
    \small
    \setlength{\tabcolsep}{4pt}
    \begin{tabular}{@{}p{4.6cm}p{6.2cm}p{2.2cm}@{}}
        \toprule
        \textbf{Path} & \textbf{Contents} & \textbf{Format} \\
        \midrule
        \texttt{listmode/reference/}          & Full-count reference list-mode stream & \todo{vendor list-mode / .l} \\
        \texttt{listmode/decayed/}            & Decayed reduced-activity list-mode, per level & \todo{vendor list-mode} \\
        \texttt{sinograms/}                   & Corrected sinograms per series (reference, decayed, thinned) & \todo{format, e.g. .s / HDF5} \\
        \texttt{recon/\{reference,decayed,thinned\}/} & Reconstructed volumes per series and realization & DICOM + NIfTI \\
        \texttt{metadata/metadata.json}       & Per-series acquisition + reconstruction parameters, count levels, seeds & JSON \\
        \texttt{ground\_truth.json}           & Sphere diameters, fill concentrations, true contrasts, NEMA ROI geometry & JSON \\
        \texttt{thinning/seed\_recipe.json}   & Deterministic seed recipe regenerating every thinned realization & JSON \\
        \texttt{manifest.sha256}              & SHA-256 of every released artifact (the content address) & text \\
        \bottomrule
    \end{tabular}
    \caption{PWM-PET-IQ 1.0 release contents. Reconstructed volumes are provided in
    both DICOM (for clinical-toolchain compatibility) and NIfTI (for analysis).
    List-mode streams are released where the scanner vendor's list-mode format
    license permits redistribution; where it does not, the thinning
    \texttt{seed\_recipe.json} plus the deposited sinograms still allow every thinned
    realization to be regenerated (\todo{confirm list-mode redistribution license}).}
    \label{tab:data_records}
\end{table}

Per-series acquisition and reconstruction parameters are summarized in
Table~\ref{tab:acq_params}; exact per-series values are in \texttt{metadata.json}.

\begin{table}[ht]
    \centering
    \small
    \setlength{\tabcolsep}{4pt}
    \begin{tabular}{lcccc}
        \toprule
        \textbf{Parameter} & \textbf{Reference} & \textbf{Decayed} & \textbf{Thinned} & \textbf{Source} \\
        \midrule
        Count level $r$                & 1.0            & \todo{levels}  & \todo{levels}  & derived \\
        True total counts (prompts)    & \todo{counts}  & \todo{counts}  & \todo{counts}  & list-mode \\
        Realizations                   & 1              & 1 per level    & \todo{N} per level & --- \\
        Frame duration (s)             & \todo{}        & \todo{}        & n/a (thinned)  & \texttt{metadata.json} \\
        Decay-corrected activity (MBq) & \todo{}        & \todo{}        & n/a            & dose calibrator \\
        Randoms fraction (\%)          & \todo{}        & \todo{}        & \todo{}        & scanner \\
        \bottomrule
    \end{tabular}
    \caption{Per-series-class acquisition parameters, extracted from the deposited
    \texttt{metadata.json}. Every cell is \emph{provisional} until the acquisition and
    the committed extraction script land (SUBMISSION\_CHECKLIST \S2--\S3); no value is
    asserted from a committed artifact in this skeleton.
    \todo{Regenerate all cells from \texttt{metadata.json} via the committed extraction
    script once acquisition completes.}}
    \label{tab:acq_params}
\end{table}

Representative reconstructed slices across the count levels are shown in
Figure~\ref{fig:examples}.

\todoFigure{fig:examples}{\textbf{Representative PWM-PET-IQ slices.} Axial slice
through the sphere plane of the NEMA NU-2 IQ phantom at each count level (full,
decayed-$0.50$, decayed-$0.25$, decayed-$0.10$, and one matched thinned realization
per level), all rendered at identical windowing, showing the count-dependent noise
increase against the fixed known geometry. Populate from the reconstructed volumes at
submission.}

%% =====================================================================
\section*{Technical Validation}

All endpoints below are computed against the physical ground-truth table
(\texttt{ground\_truth.json}); none requires a reader panel. Every number is emitted
by a committed analysis script over the deposited records, and no endpoint is asserted
in this skeleton.

\subsection*{Contrast recovery vs count level}
For each sphere and each count level we compute the NEMA NU-2 contrast-recovery
coefficient (CRC)~\citep{nema_nu2} against the known true contrast, on the reference,
decayed, and thinned series. This is the primary image-quality endpoint and
establishes how contrast recovery degrades with decreasing counts for a fixed
reconstruction.

\todoTable{tab:crc}{Contrast-recovery coefficient per sphere (10--37\,mm) at each
count level, for the reference, decayed, and thinned series, against physical ground
truth. Primary image-quality endpoint.}

\subsection*{Background variability and contrast-to-noise}
Per NEMA NU-2, we compute background variability from the 12-region ROI set and the
per-sphere contrast-to-noise ratio at each count level, on decayed and thinned series,
so users can read the noise--count tradeoff directly.

\todoFigure{fig:cnr_vs_count}{Background variability and per-sphere contrast-to-noise
ratio as a function of count level $r$, for decayed (physical) and thinned (list-mode)
series. Quantifies the noise--count tradeoff at fixed reconstruction.}

\subsection*{Decayed-vs-thinned agreement (the key validation)}
Because the decayed and thinned series share phantom geometry exactly, we can measure
how faithfully list-mode Poisson thinning reproduces a genuinely reduced-activity
acquisition. At each matched count level we report the per-voxel and per-ROI agreement
(bias and limits of agreement) between the decayed reference and the thinned
realizations, and the difference in randoms fraction and noise texture. This
quantifies --- rather than assumes --- the validity of the thinning model that
underlies most low-count PET simulators and the PET signal-reduction operator of the
companion framework~\citep{ws2framework}.

\todoFigure{fig:decayed_vs_thinned}{Per-level agreement between physically-decayed and
list-mode-thinned reduced-count data (Bland--Altman on sphere CRC and background
variability; randoms-fraction and noise-texture deltas). The central validation
figure: it measures the thinning-model error.}

\subsection*{Realization reproducibility}
For the thinned series we confirm that the released seed recipe regenerates every
realization bit-for-bit from the reference list-mode (or sinogram) plus the published
seeds, and we report the across-realization variance at each count level as the
noise-replicate budget available to downstream denoiser evaluation.

\todoTable{tab:reproducibility}{Across-realization mean and variance of sphere CRC and
background variability at each thinned count level, and confirmation of bit-for-bit
seed regeneration. Establishes the noise-replicate budget.}

%% =====================================================================
\section*{Usage Notes}

\subsection*{Loading the data}

Reconstructed volumes are released as both DICOM and NIfTI and can be read with any
standard toolchain; the harmonized per-series metadata and ground-truth table are
plain JSON. A thin loader exposes series, count level, realization index, and the
ground-truth contrasts under one interface:

\begin{verbatim}
    from pwm_pet_iq import PetIqDataset
    ds = PetIqDataset(root="/path/to/pwm_pet_iq_1_0/")
    s  = ds.series(kind="thinned", r=0.25, realization=0)
    # s.volume        -> ndarray [Z, Y, X] (recon)
    # s.count_level   -> 0.25
    # s.ground_truth  -> per-sphere true diameter / concentration / contrast
    # s.metadata      -> acquisition + reconstruction parameters
\end{verbatim}

\subsection*{Analyses uniquely enabled by 1.0}

\begin{enumerate}
    \item \textbf{Ground-truth-anchored low-count reconstruction benchmarking.}
    Score a reduced-activity reconstruction method against \emph{physical} sphere
    contrast, not against a full-dose reconstruction of unknown error.
    \item \textbf{Decayed-vs-thinned model validation.} Test whether a method
    developed on list-mode-thinned data transfers to genuinely reduced-activity data
    of the same geometry --- a check no single-origin phantom release supports.
    \item \textbf{Noise-replicate denoiser evaluation.} Use the multiple seeded
    thinned realizations per level to estimate the variance of a denoiser's output,
    not just a point estimate.
\end{enumerate}

\subsection*{Consuming the dataset with the signal-equivalence framework}

The companion \texttt{pwm\_dose\_equivalence} tool~\citep{ws2framework} reads this
dataset directly: with the phantom's known contrasts as ground truth and the
contrast-recovery coefficient as the task metric, a single call issues a
reduced-activity signal-equivalence credential for any reconstruction method against
the full-count reference. This is an \emph{optional} consumer of the dataset; the
dataset stands alone and its Technical Validation does not depend on the framework.

\subsection*{Limitations}

\paragraph{Phantom, not patient.} The NEMA NU-2 IQ phantom captures contrast recovery
and background noise for spherical lesions on a uniform background; it does not
reproduce patient anatomy, respiratory motion, or heterogeneous tracer uptake.
Conclusions about clinical low-count performance must be confirmed on patient data.

\paragraph{Single scanner / single tracer in 1.0.} The 1.0 release is acquired on one
scanner with one tracer at one sphere-to-background ratio (\todo{SBR}). Multi-scanner,
multi-tracer, and multi-SBR extensions are deferred to a future release and are out of
1.0 scope.

\paragraph{Thinning omits activity-dependent detector physics.} List-mode thinning
matches coincidence-count statistics but not dead-time, pile-up, or the
activity-dependent randoms fraction; the decayed series are provided precisely so that
this gap is measured (Technical Validation) rather than hidden.

\paragraph{Reliance on continued hosting.} The SHA-256 manifest records what 1.0
\emph{is} independent of any host, but does not guarantee the deposited files remain
downloadable if a repository withdraws them; the PhysioNet DOI and Zenodo mirror are
provided to reduce this risk.

%% =====================================================================
\subsection*{Data versioning and persistent identifiers}

Each release is identified by the SHA-256 hash of \texttt{manifest.sha256}, which
hashes every deposited artifact. Releases follow \texttt{MAJOR.MINOR.PATCH} semantics:
MAJOR breaks the metadata/ground-truth schema, MINOR adds scanners/tracers/SBR points,
and PATCH corrects errors. The primary citation handle is the PhysioNet DOI; the Zenodo
record mirrors the same content hash.

%% =====================================================================
\section*{Code Availability}

\begin{itemize}
    \item Loader: \url{https://pypi.org/project/pwm_pet_iq/} (\todo{license}) --- \todo{publish to PyPI}
    \item List-mode thinning + reconstruction container:
    \url{https://github.com/integritynoble/low_dose_CT/tree/main/WS-2b_pet_phantom/pipelines/} --- \todo{add pipelines}
    \item Analysis scripts regenerating every Technical-Validation table/figure:
    \url{https://github.com/integritynoble/low_dose_CT/tree/main/WS-2b_pet_phantom/analysis/} --- \todo{add analysis scripts}
    \item Deposited records + manifest: PhysioNet DOI \todo{DOI}, mirrored at Zenodo \todo{Zenodo DOI}.
\end{itemize}

%% =====================================================================
\section*{Author contributions}
\todo{Fill at submission per CRediT taxonomy.}

\section*{Competing interests}
\todo{Declare at submission. Note PWM Protocol Foundation affiliation and state absence
of financial interest in any reconstruction method benchmarked on this dataset.}

\section*{Acknowledgments}
\todo{Fill at submission. Acknowledge the PET physics / radiochemistry staff who
prepared and acquired the phantom, and the open-source maintainers of the
reconstruction and analysis toolchain.}

%% =====================================================================
\bibliographystyle{plainnat}
\bibliography{refs}

\end{document}
